\chapter{Conclusions}
\label{ch:conclusions}

\section{Summary}
\label{sec:summary}

This thesis set out to answer a practical engineering question: what is
the smallest, most maintainable offline-first architecture sufficient for
a CRUD-style field-reporting mobile application, and what trade-offs does
that architecture force a developer to make explicit? Using a campus
facility-maintenance reporting scenario (CampusFix) as a concrete
instance of the broader field-data-collection problem
(\cref{ch:introduction}), the thesis:

\begin{itemize}
  \item derived a requirements set (\cref{ch:requirements}) that
  separates ordinary functional requirements from eight offline-specific
  requirements (OFR-1--OFR-8) covering local-first writes, offline session
  continuity, deterministic local identity, explicit sync state,
  connectivity-triggered synchronization, resumability, conflict
  resolution, and no-data-loss-on-failure;
  \item designed a four-layer architecture (\cref{ch:design}) in which the
  presentation layer depends only on a local Room/SQLite database, with a
  dedicated synchronization layer bridging that database to a Firebase
  Firestore backend, and a whole-record last-write-wins conflict policy
  chosen for its simplicity and adequacy for single-owner records;
  \item implemented that architecture (\cref{ch:implementation}) as a
  Kotlin Multiplatform application sharing 82.6\% of its Kotlin source
  (2{,}326 of 2{,}817 lines) between Android and iOS, with
  platform-specific code confined to connectivity monitoring and thin
  entry-point glue; and
  \item evaluated the resulting system (\cref{ch:evaluation}) against six
  scenarios -- offline creation, delayed sync, failed sync, large sync
  queues, conflicting edits, and end-user usability -- deriving several
  conclusions directly from the shape of the implemented algorithm rather
  than only from measurement, including a precise, code-traced account of
  how the system currently behaves under the conflicting-edit scenario.
\end{itemize}

The overall finding is that a relatively small, well-isolated
synchronization module -- in this implementation, a single
\texttt{SyncManager} class of under 100 lines plus a platform-abstracted
\texttt{ConnectivityObserver} -- is sufficient to make a CRUD-style field
reporting application fully usable offline, provided the architectural
discipline of \cref{sec:architecture-overview} (UI depends only on local
storage, never on the network) is maintained throughout. Getting the
\emph{local-first} half of offline-first right is, in this experience,
largely a matter of architectural discipline rather than algorithmic
difficulty; getting the \emph{synchronization} half right -- specifically,
bidirectional conflict resolution -- is where genuine algorithmic and
engineering effort remains, as discussed next.

\section{Limitations}
\label{sec:limitations}

The following limitations of the current system should be read alongside
its contributions.

\begin{enumerate}
  \item \textbf{Synchronization is one-directional (push-only).} As
  established in \cref{sec:conflict-design} and demonstrated by tracing
  the code in \cref{sec:conflict-eval}, \texttt{SyncManager} uploads local
  changes to Firestore but never reads the current remote document before
  overwriting it, and never pulls remote-only changes down to the device.
  The designed bidirectional last-write-wins policy (OFR-7) is therefore
  only partially realized. This is the single most consequential
  limitation identified in this thesis.
  \item \textbf{No explicit \texttt{FAILED} transition on upload error.}
  \cref{fig:sync-state-machine} reserves a \texttt{FAILED} status for a
  report whose upload has been attempted and rejected, but the current
  \texttt{uploadPending} implementation (\cref{lst:syncmanager}) simply
  leaves a report in its prior state (\texttt{PENDING}) on exception
  rather than explicitly marking it \texttt{FAILED}; the user therefore
  cannot currently distinguish "not yet attempted" from "attempted and
  failed" purely from the status shown in the UI.
  \item \textbf{Sequential, unbatched uploads.} \cref{sec:eval-large-queue}
  identifies that sync time scales linearly with queue size because
  reports are uploaded one at a time in a single coroutine, with no
  batching or parallelism.
  \item \textbf{No server-side authorization.} \cref{sec:constraints}
  notes that the \texttt{owner} field is a client-set, trusted value used
  for filtering, not enforced by Firestore security rules in the current
  implementation; a malicious or buggy client could, in principle, read or
  write another user's reports.
  \item \textbf{Photo synchronization is designed but not fully wired.}
  The schema (\cref{tab:report-fields}) reserves a
  \texttt{remote\_photo\_url} field for a Firebase Storage upload (FR-13),
  but the evaluated \texttt{SyncManager} implementation
  (\cref{lst:syncmanager}) only synchronizes the report's structured
  fields, not the photo binary itself.
  \item \textbf{No periodic background sync.} As discussed in
  \cref{sec:sync-triggers}, synchronization is triggered by connectivity
  changes, app start, and manual action, but not by a periodic background
  job, so a device that regains connectivity while the app is fully
  backgrounded/terminated will not synchronize until the app is next
  opened.
\end{enumerate}

\section{Future Work}
\label{sec:future-work}

Ordered roughly by priority, given the limitations above:

\begin{enumerate}[label=\textbf{F\arabic*.}]
  \item \textbf{Bidirectional synchronization with real last-write-wins.}
  Extend \texttt{SyncManager} to, before every \texttt{set()}, read the
  current remote document and compare its \texttt{updatedAt} against the
  local report's \texttt{updatedAt}, keeping whichever is newer; and add a
  pull path (either a one-shot \texttt{get()} per sync or a Firestore
  snapshot listener) so that remote-only changes (e.g. a report edited
  through a hypothetical facilities-management back office) reach the
  device. A Firestore transaction, rather than a plain read-then-write,
  would additionally close the small race window between the read and the
  write.
  \item \textbf{Explicit \texttt{FAILED} status on upload error.} A
  small, low-risk change to \texttt{uploadPending}'s \texttt{catch} block
  to call \texttt{dao.update(report.copy(syncStatus = SyncStatus.FAILED))},
  closing limitation 2 above and making the "failed sync" scenario of
  \cref{sec:eval-failed-sync} fully observable from the UI.
  \item \textbf{Batched/parallel uploads.} Replace the sequential
  \texttt{for} loop in \texttt{uploadPending} with either a bounded-
  concurrency \texttt{coroutineScope} that uploads several reports at
  once, or Firestore's native batched-write API, to flatten the linear
  scaling observed in \cref{sec:eval-large-queue}.
  \item \textbf{Firestore security rules keyed on Firebase Auth UID.}
  Replace the trusted, client-set \texttt{owner} email field with
  server-enforced rules that only allow a document's true creator (by
  Firebase Auth UID) to write it, closing limitation 4.
  \item \textbf{Complete photo synchronization.} Wire the existing
  \texttt{remotePhotoUrl} field to an actual Firebase Storage upload,
  triggered alongside the report's own document sync.
  \item \textbf{Periodic background sync.} Add a platform background task
  (\texttt{WorkManager} periodic work on Android; a
  \texttt{BGAppRefreshTask} on iOS) as a fourth synchronization trigger,
  for the case where connectivity returns while the app is not in the
  foreground.
  \item \textbf{Field-level conflict resolution.} Once bidirectional sync
  (F1) is in place, investigate whether whole-record LWW's
  data-loss failure mode (\cref{fig:conflict-scenario}) is frequent enough
  in real usage to justify the added complexity of field-level merge or a
  CRDT-based representation~\cite{shapiro2011crdt}.
  \item \textbf{Automated tests for the synchronization module.} The
  evaluation in \cref{ch:evaluation} is manual/scenario-based; an
  automated test suite exercising \texttt{SyncManager} against a Firestore
  emulator (including the large-queue and conflict scenarios) would let
  these properties be checked on every code change rather than only during
  a dedicated evaluation pass.
  \item \textbf{Repeat the evaluation of \cref{ch:evaluation} with real
  measurements}, replacing every table and figure marked
  "illustrative" with data from actual test runs, and re-examining the
  discussion in \cref{sec:eval-discussion} against those numbers.
\end{enumerate}

\section{Final Remarks}
\label{sec:final-remarks}

Offline-first design is often presented, in industry writing, as either a
trivial addition of "a local cache" or as an unavoidably hard problem
requiring CRDTs and vector clocks. The experience of building CampusFix
suggests the practical reality, at least for CRUD-style field-reporting
applications, sits between those two extremes: making an application
\emph{fully usable offline} is genuinely tractable with a disciplined
layered architecture, a local database as the single source of truth, and
a connectivity-aware sync module -- the kind of code shown in
\cref{ch:implementation} is neither exotic nor large. Making
synchronization \emph{fully correct} under concurrent, disconnected edits
is the part that remains hard, and where this thesis's honest account of
its own implementation gap (\cref{sec:conflict-eval}) is, perhaps, its
most useful methodological contribution: a working offline-first
application and a \emph{correctly conflict-free} offline-first
application are not the same claim, and a thesis -- like a codebase --
should be explicit about which one it is making.
